\subsection{群上函数的卷积}

我们保留 No.~4 的记号 G、\(\beta\)、\(\chi\)、\(\chi'\)。

回忆：若 f 是 G 上的复值函数，则局部 \(\beta\)-可积这一性质与 \(\beta\) 的选取无关。设 \(\mathcal{L}(\mathrm{G})\) 为具有这一性质的函数组成的集合。若 \(f \in \mathcal{L}(\mathrm{G})\)，\(g \in \mathcal{L}(\mathrm{G})\)，则关系

\[\text{«}f \cdot \beta\text{ 与 }g \cdot \beta\text{ 可作卷积 »}\]

与 \(\beta\) 的选取无关（§3, No.~1 命题 6）。这时我们就说 \(f\) 与 \(g\) 可作卷积。由 No.~1，\((f\cdot \beta)*(g\cdot \beta)\) 具有 \(h\cdot \beta\) 的形式，其中 \(h\in \mathcal{L}(\mathbf{G})\)，\(h\) 在局部 \(\beta\)-可忽略集的范围内被确定。我们将记 \(h = f * ^ {\beta} g\)，并说 \(h\) 是 \(f\) 与 \(g\) 相对于 \(\beta\) 的一个卷积。（当不致混淆时省略 \(\beta\)。）若把 \(\beta\) 换成 \(\psi \cdot \beta\)，其中 \(\psi\) 是 \(\mathbf{G}\) 在 \(\mathbf{R}_+^*\) 中的一个连续表示，则 \(h\) 不变（§3, No.~1 命题 6）；若把 \(\beta\) 换成 \(a\beta\)（\(a\in \mathbf{R}_+^*\)），则 \(h\) 换成 ah。G 上若干个函数的卷积以类似方式定义。

若 f 与 g 的卷积之一连续，则它被唯一确定，因为 \(\beta\) 的支集是 G。这时称它为 f 与 g 相对于 \(\beta\) 的卷积。

显然

\[
f * ^ {\beta} g = (f \cdot \beta) * ^ {\beta} g = f * ^ {\beta} (g \cdot \beta).\tag{14}
\]

命题 9. --- 设 \(f,g\) 属于 \(\mathcal{L}(\mathbf{G})\)。假设函数 \(s\mapsto g(s^{-1}x)f(s)\chi(s^{-1})\) 是本质 \(\beta\)-可积的，除去一个局部 \(\beta\)-可忽略的 \(x\) 值集合，并且函数

\[
x \mapsto \int | g (s ^ {- 1} x) f (s) | \chi (s ^ {- 1}) d \beta (s),
\]

它局部 \(\beta\)-几乎处处有定义且局部 \(\beta\)-可积。则 \(f\) 与 \(g\) 可作卷积。

这由 No.~1 的命题 1 得出。

命题 10. --- 设 \(f, g\) 属于 \(\mathcal{L}(G)\)。假设这两个函数之一连续，或在一列紧集的并的余集上为零。若 \(f\) 与 \(g\) 可作卷积，则函数 \(f * g\) 在局部 \(\beta\)-几乎处处由下式给出

\[
\begin{array}{c} (f * g) (x) = \int_ {\mathrm{G}} g (s ^ {- 1} x) f (s) \chi (s ^ {- 1}) d \beta (s) \\ = \int_ {\mathrm{G}} f (x s ^ {- 1}) g (s) \chi^ {\prime} (s ^ {- 1}) d \beta (s). \end{array}\tag{15}
\]

这由 No.~1 的命题 2 以及 No.~4 中的备注得出。

特别地，若 \(f * g\) 连续且在每一点由 (15) 给出，则

\[
(f * g) (e) = \int g (s ^ {- 1}) f (s) \chi (s ^ {- 1}) d \beta (s) = \int f (s ^ {- 1}) g (s) \chi^ {\prime} (s ^ {- 1}) d \beta (s).\tag{16}
\]

更特别地，若 \(\beta\) 是左 Haar 测度且是右 Haar 测度，且 \(f*g\) 与 \(g*f\) 连续并在每一点由 (15) 与 \(g*f\) 的类似公式给出，则

\[
(f * g) (e) = (g * f) (e) = \int f (s) g (s ^ {- 1}) d \beta (s).\tag{17}
\]

命题 11. --- 设 \(f,g\) 属于 \(\mathcal{L}(\mathrm{G})\)。假设函数 \(f,g\) 之一连续，且函数 \(f,g\) 之一具有紧支集。则 \(f\) 与 \(g\) 可作卷积。公式 (15) 对所有 \(x\in\mathrm{G}\) 定义了一个连续的乘积 \(f*g\)。若 \(f\in\mathcal{K}(\mathrm{G})\) 且 \(g\in\mathcal{K}(\mathrm{G})\)，则 \(f*g\in\mathcal{K}(\mathrm{G})\)。

这由 No.~2 的命题 3 与命题 4 得出。

命题 12. --- 设 \(p\) 与 \(q\) 为两个共轭指数 \((1 \leqslant p \leqslant +\infty)\)。若 \(f\chi^{-1/q} \in \mathrm{L}^1(\mathrm{G}, \beta)\) 且 \(g \in \mathrm{L}^p(\mathrm{G}, \beta)\)，则 \(f\) 与 \(g\) 可作卷积，\(f * g\) 在局部 \(\beta\)-几乎处处等于 \(\mathrm{L}^p(\mathrm{G}, \beta)\) 中的一个函数，且

\[
\left\| f * g \right\| _ {p} \leqslant \left\| f \chi^ {- 1 / q} \right\| _ {1} \left\| g \right\| _ {p}.
\]

若 \(f \in \mathrm{L}^p(\mathbf{G},\beta)\) 且 \(g\chi'^{-1/q} \in \mathrm{L}^1(\mathbf{G},\beta)\)，则 \(f\) 与 \(g\) 可作卷积，\(f * g\) 在局部 \(\beta\)-几乎处处等于 \(\mathrm{L}^p(\mathbf{G},\beta)\) 中的一个函数，且

\[
\| f * g \| _ {p} \leqslant \| f \| _ {p} \| g \chi^ {\prime - 1 / q} \| _ {1}.
\]

这由 No.~2 的命题 5 与命题 6 以及 No.~4 中的备注得出。

命题 13. --- 若 \(f\chi^{-1} \in \mathbf{L}^1(\mathbf{G}, \beta)\) 且 \(g \in \overline{\mathcal{H}(\mathbf{G})}\)，或者若 \(f \in \overline{\mathcal{H}(\mathbf{G})}\) 且 \(g\chi'^{-1} \in \mathbf{L}^1(\mathbf{G}, \beta)\)，则 \(f\) 与 \(g\) 可作卷积，且 (15) 对每一 \(x \in \mathbf{G}\) 定义了一个乘积 \(f * g\)，它属于 \(\overline{\mathcal{H}(\mathbf{G})}\)。

这由 No.~2 的命题 5 以及 No.~4 中的备注得出。

命题 14. --- 若 \(f\chi^{-1} \in \mathbf{L}^1(\mathbf{G}, \beta)\) 且 \(g \in \mathbf{L}^\infty(\mathbf{G}, \beta)\)，则公式 (15) 对每一 \(x \in \mathbf{G}\) 定义了一个有界的乘积 \(f * g\)，它对 \(\mathbf{G}\) 的右一致结构是一致连续的。

我们已经知道 \(f * g\) 属于 \(\mathrm{L}^{\infty}(\mathrm{G},\beta)\)（No.~2 命题 5）；此外，\((f*g)(x) = \int f(xs^{-1})g(s)d\nu(s)\)，其中置 \(\nu = \chi'^{-1}\cdot \beta\)；\(\nu\) 是一个右 Haar 测度。因此

\[
\begin{array}{c} | (f * g) (x) - (f * g) (x ^ {\prime}) | \leqslant \| g \| _ {\infty} \int | f (x s ^ {- 1}) - f (x ^ {\prime} s ^ {- 1}) | d \nu (s) \\ = \| g \| _ {\infty} \int | (f (s ^ {- 1}) - f (x ^ {\prime} x ^ {- 1} s ^ {- 1}) d \nu (s) \end{array}
\]

且后一积分可以任意小，只要 \(x'x^{-1}\) 属于 \(e\) 的一个适当邻域（§2, No.~5 命题 8）。

命题 15. --- 设 \(p\) 与 \(q\) 为两个共轭指数 \((1 < p < +\infty)\)。假设 \(\beta\) 是左不变的。设 \(f \in \mathrm{L}^p(\mathrm{G}, \beta)\)，\(g \in \mathrm{L}^q(\mathrm{G}, \breve{\beta})\)。则 \(f\) 与 \(g\) 可作卷积。公式 (15) 对每一 \(x \in \mathrm{G}\) 定义了一个乘积 \(f * g\)，它属于 \(\overline{\mathcal{H}(\mathrm{G})}\) 且满足 \(\| f * g \|_{\infty} \leqslant \| f \|_p \| \breve{g} \|_q\)。

因为，我们有 \(\check{g} \in \mathbf{L}^q(\mathbf{G}, \beta)\)，因此函数 \(s \mapsto g(s^{-1}x)f(s)\) 是 \(\beta\)-可积的，对每一 \(x \in \mathbf{G}\)。此外，

\[
\begin{array}{c} \int | g (s ^ {- 1} x) f (s) | d \beta (s) \leqslant \bigg (\int | f (s) | ^ {p} d \beta (s) \bigg) ^ {1 / p} \bigg (\int | g (s ^ {- 1} x) | ^ {q} d \beta (s) \bigg) ^ {1 / q} \\ = \| f \| _ {p} \bigg (\int | \check {g} (x ^ {- 1} s) | ^ {q} d \beta (s) \bigg) ^ {1 / q} = \| f \| _ {p} \| \check {g} \| _ {q}, \end{array}
\]

因此 \(f\) 与 \(g\) 可作卷积（命题 9）。同时可以看到 (15) 对每一 \(x\) 定义了一个乘积 \(f * g\)，使得

\[
\left| (f * g) (x) \right| \leqslant \| f \| _ {p} \| \check {g} \| _ {q}.
\]

对 \(f, g\)（\(\mathcal{H}(\mathrm{G})\) 中的），我们有 \(f * g \in \mathcal{H}(\mathrm{G})\)（命题 11）；因此，对 \(f \in \mathrm{L}^p(\mathrm{G}, \beta)\) 与 \(g \in \mathrm{L}^q(\mathrm{G}, \beta)\)，由 (15) 给出的乘积 \(f * g\) 是 \(\mathcal{H}(\mathrm{G})\) 中函数的一致极限，故属于 \(\overline{\mathcal{H}(\mathrm{G})}\)。

推论. --- 设 \(f \in \mathbf{L}^2(\mathbf{G}, \beta)\)，\(g \in \mathbf{L}^2(\mathbf{G}, \beta)\)。则 \(f\) 与 \(\breve{g}\) 可作卷积。卷积 \(f * \breve{g}\) 之一属于 \(\overline{\mathcal{H}(\mathbf{G})}\)，且它在 \(e\) 处的值为 \(\int_{\mathbf{G}} f(s) \overline{g(s)} d\beta(s)\)。

只需在命题 15 中取 \(p = q = 2\) 并应用 (16)。

我们不再假设 \(\beta\) 是左不变的。设 \(\rho\) 为 G 上的一个下半连续的有限函数 \textgreater0，使得对 G 中所有 s, t 有 \(\rho(st) \leqslant \rho(s)\rho(t)\)。我们用 \(\mathrm{L}^{\rho}(\mathrm{G}, \beta)\) 表示 G 上对 \(\rho \cdot \beta\) 可积的复值函数的等价类组成的集合。借助映射 \(f \mapsto f \cdot \beta\)，\(\mathrm{L}^{\rho}(\mathrm{G}, \beta)\) 可等同于 \(\mathcal{M}^{\rho}(\mathrm{G})\) 中以 \(\beta\) 为基的元素组成的集合（该集合与 \(\beta\) 的选取无关）。若置

\[
\| f \| _ {\rho} = \int_ {G} | f (s) | \rho (s) d \beta (s)
\]

对 \(f \in \mathrm{L}^{\rho}(\mathrm{G},\beta)\)，这一等同与范数相容，于是 \(\mathrm{L}^{\rho}(\mathrm{G},\beta)\) 表现为 \(\mathcal{M}^{\rho}(\mathrm{G})\) 的一个完备赋范子代数。由 §3, No.~2 命题 10，它甚至是 \(\mathcal{M}^{\rho}(\mathrm{G})\) 的一个双边理想。（取 \(\rho = 1\)，就重新得到 No.~4 的断言之一。）特别地，\(\mathrm{L}^{1}(\mathrm{G},\beta)\) 可等同于 \(\mathcal{M}^{1}(\mathrm{G})\) 的一个闭双边理想。

命题 16. --- 设 U 为 G 在 Banach 空间 E 中的一个连续表示。置 \(\rho(s) = \|U(s)\|\)（对所有 \(s \in G\)）。对每一 \(f \in L^{\rho}(G, \beta)\)，置 \(U(f) = U(f \cdot \beta)\)。则 \(f \mapsto U(f)\) 是代数 \(L^{\rho}(G, \beta)\) 在 E 中的一个线性表示，满足 \(\|U(f)\| \leqslant \|f\|_{\rho}\)。

这由 §2, No.~6 与 §3, No.~3 命题 11 得出。

